% ===== 原文 pp.26–28 =====
% 译者 tr-04：§4 开头、§4.1、§4.2 前半（含表 4.7、表 4.8 之前的铺垫）

% [原文 p.26]
\section{Higgs 丛模空间中的膜}

\begin{quote}
\itshape 如果人们不相信数学是简单的，那只是因为他们还没有意识到生活有多么复杂。

\normalfont\begin{flushright}
John von Neumann
\end{flushright}
\end{quote}

Higgs 丛（以及平坦联络）在弦论与几何 Langlands 纲领中的出现，促使研究者去研究这些模空间的凝聚层的导出范畴（derived category of coherent sheaves）与 Fukaya 范畴。因此，理解 Higgs 丛模空间中支撑全纯层的拉格朗日子流形（A-膜）以及它们的对偶对象（B-膜）已成为一项基本任务。Hitchin 方程 (2.7)--(2.8) 的解空间是一个超凯勒流形，因此存在一族复结构；遵循 [17, 60, 69] 的记号，我们在其中固定满足四元数关系的 $I$、$J$、$K$。特别地，在此约定下，$\cM_{G_\C}$ 的光滑轨迹对应于带有结构 $I$ 的 Hitchin 方程的解空间。在本讲义中我们将采用物理学家的说法：支撑平坦联络的拉格朗日子流形称为 A-膜（A-brane），支撑复层的复子流形称为 B-膜（B-brane）。通过考察膜的支撑，我们会说超凯勒流形的子流形相对于其中每个复结构分别是 A 型或 B 型的，因此可以考虑如下类型的膜
\[
(\mathrm{B},\mathrm{B},\mathrm{B}),\quad (\mathrm{B},\mathrm{A},\mathrm{A}),\quad (\mathrm{A},\mathrm{B},\mathrm{A})\quad\text{以及}\quad(\mathrm{A},\mathrm{A},\mathrm{B}).
\]

% [Tutor Note: p.26 首段把"支撑全纯层"的拉格朗日子流形称为 A-膜，与本页下文 A-膜支撑平坦联络、B-膜支撑全纯层的定义相反，疑为原文笔误；照原样译出。]

% [Tutor Note: p.26 膜类型三元组的约定：$(\mathrm{B},\mathrm{B},\mathrm{B})$ 这类三元组的分量依次对应复结构 $I,J,K$，"B"表示关于该复结构是复子流形（支撑复层），"A"表示关于该复结构是拉格朗日子流形（支撑平坦联络）。可用 p.30 的 $i_1,i_2,i_3$ 核对：$i_1$ 固定 $(\mathrm{B},\mathrm{A},\mathrm{A})$ 型、$i_2$ 固定 $(\mathrm{A},\mathrm{B},\mathrm{A})$ 型、$i_3$ 固定 $(\mathrm{A},\mathrm{A},\mathrm{B})$ 型。]

本章我们将通过在黎曼面 $\Sigma$ 上与复李群 $G_\C$ 上增添附加结构，来构造、描述并研究这类膜：

\begin{itemize}[leftmargin=1.4em,itemsep=1pt,topsep=1pt]
\item \textbf{通过黎曼面 $\Sigma$ 上有限群作用的膜。}有限群作用 $\Gamma\times\Sigma\to\Sigma$ 在紧连通黎曼面上早已被研究；通过这些作用，我们将按照 [22, 59] 考虑经由 $\Gamma$-等变表示出现的 $(\mathrm{B},\mathrm{B},\mathrm{B})$-膜。
\item \textbf{通过反全纯对合的膜。}我们此前已经看到，反全纯对合可用于定义实 Higgs 丛，这一想法始于 [62]。通过在黎曼面上考虑实结构以及它们与群对合的复合，我们将按照 [15, 17]（另见 [90]）定义 $(\mathrm{B},\mathrm{A},\mathrm{A})$、$(\mathrm{A},\mathrm{B},\mathrm{A})$ 与 $(\mathrm{A},\mathrm{A},\mathrm{B})$-膜。在最后一章按照 [85] 考虑超多边形时，这些构造将扮演特殊的角色。
\end{itemize}

\subsection{通过有限群作用的膜}

正如 [59] 所示，在复 Higgs 丛的模空间中，定义 $(\mathrm{B},\mathrm{B},\mathrm{B})$-膜的一个自然途径是通过有限群 $\Gamma$ 在亏格为 $g$ 的底黎曼面上的作用。当亏格为 2 或 3 时，所有有限群作用的完全分类出现在 [31, Tables 4, 5] 中，这使人们能够在这些情形下进行显式计算。此外，对于经由 $\Sigma$ 的自同构在秩 2 丛上诱导出的作用，[3] 用抛物结构给出了不动点的一个非常具体的描述。

% [Tutor Note: p.26 §4.1 首句原文 "is though the action of a finite group" 中 "though" 应为 "through"；另 p.26 首段 "Throughout this notes" 应为 "Throughout these notes"。两处均照原样理解并译出。]

% [原文 p.27]
在此框架下，一个被大量研究的问题是 $\Sigma$ 上平坦 $\Gamma$-等变 $G_\C$-联络的出现，为此需要固定底向量丛的一个 $C^\infty$ 平凡化 $\underline{\C}^n=\Sigma\times\C^n\to\Sigma$。关于黎曼面 $\Sigma$ 上的 $\Gamma$-等变平坦联络有几种不同的观点；这里我们采用 [59] 的定义：$\Gamma$-等变平坦联络是指 $\underline{\C}^n\to\Sigma$ 上的一个平坦联络 $\nabla$，使得对每个 $\phi\in\Gamma$ 都存在 $\Gl(n,\C)$-规范变换 $g_\phi\colon\Sigma\to\Gl(n,\C)$，满足
\begin{itemize}[leftmargin=1.4em,itemsep=1pt,topsep=1pt]
\item $\phi^*\nabla=\nabla\cdot g_\phi$，
\item 且 $\phi\mapsto g_\phi$ 是\textit{广义群同态}（generalized group homomorphism），即满足
\[
\begin{gathered}
g_{\id}=\id,\\
g_{(\phi\circ\tau)}(p)=g_{\tau(p)}\circ g_\phi(\tau(p)).
\end{gathered}
\]
\end{itemize}

% [Tutor Note: p.27 原文 "the pull-back of a flat $G_\C(n,\C)$-connection on $\Sigma/\Gamma$" 中的 $G_\C(n,\C)$ 疑为 $G_\C$ 之误（对照同段下文 $G_\C$-connection、以及例 4.1 前后的写法）；照原样排。]

取 $G_\C=\Gl(n,\C)$、$\Sl(n,\C)$，并设 $\Gamma\times\Sigma\to\Sigma$ 是由无不动点的全纯自同构给出的有限群作用。此时，任何 $\Gamma$-等变平坦 $G_\C$-联络都是 $\Sigma/\Gamma$ 上某个平坦 $G_\C(n,\C)$-联络的拉回（例如见 [59, Proposition 1]）。当作用有不动点时，不动点在 $\Sigma/\Gamma$ 中的像 $B\subset\Sigma/\Gamma$ 给出分歧覆盖
\[
\pi_\Gamma\colon\Sigma\to\Sigma/\Gamma
\]
的分支点，从而任何 $\Gamma$-等变平坦 $G_\C$-联络都是 $\Sigma/\Gamma-B$ 上某个平坦 $G_\C$-联络的拉回（例如见 [59, Proposition 2]）。由 $\Gamma$-等变平坦联络的定义可知，$\Sigma$ 上 $\Gamma$-等变平坦不可约 $G_\C$-联络 $\nabla$ 会产生一个 $\Gamma$-等变联络 $\nabla\cdot g$，其中 $g\colon\Sigma\to G_\C$ 是规范变换；它对应于 $\Sigma/\Gamma$ 上平坦（可能奇异）$G_\C$-联络的模空间中的同一个点（例如见 [59, Lemma 1]）。

\begin{exmp}{4.1}
为说明等变联络（从而等变 Higgs 丛）的出现，考虑 [59] 中如下的例子。设 $\Sigma$ 是亏格为 3 的紧黎曼面，它容许一个无不动点的对合 $\phi\colon\Sigma\to\Sigma$。通过这一作用，得到亏格为 2 的商黎曼面 $M:=\Sigma/\Z_2$ 的双重覆盖。

设 $\nabla$ 是一个非自对偶的平坦酉线丛联络，且满足 $\phi^*\nabla=\nabla^*$。则
\[
\nabla\oplus\nabla^*
\]
是 $\Sigma$ 上的 $\Z_2$-等变平坦联络，并对应 $\Sigma/\Z_2$ 上的一个不可约平坦联络。
\end{exmp}

% [Tutor Note: p.27 例 4.1 的亏格核对：由 Riemann--Hurwitz 公式，无分歧点的双重覆盖满足 $2g-2=2(2g'-2)$；代入 $g=3$ 得 $g'=2$，与商曲面 $M=\Sigma/\Z_2$ 的亏格 2 相符。]

% [原文 p.28]
通过在亏格 $g\ge 2$ 的黎曼面 $\Sigma$ 上考虑群 $\Gamma$ 的不同有限群作用，[59, Theorem 8] 证明了：不可约 $\Gamma$-等变平坦 $G_\C$-联络的规范等价类空间的连通分支，是 $\Sigma$ 上平坦不可约 $G_\C$-联络的模空间的超凯勒子流形，因此给出 $G_\C$-Higgs 丛模空间中的 $(\mathrm{B},\mathrm{B},\mathrm{B})$-膜。

当 $G_\C=\Sl(2,\C)$ 时，由上述结果，$\Sigma$ 上平坦 $\Gamma$-等变 $\Sl(2,\C)$-联络的模空间的任何连通分支，都可等同于亏格为 $\gamma$ 的 $n$-带孔紧黎曼面上的平坦 $\Sl(2,\C)$-联络或 $\PSL(2,\C)$-联络的模空间，其中 $n\in\N$ 满足某些相容性条件。此外，关于有限群在 Higgs 丛模空间上诱导出的对合，其不动点集的类型还可以说得更多：

\begin{thm}{4.2（[59] 定理 11）}
设 $\Sigma$ 是亏格 $g\ge 2$ 的紧黎曼面，$\Gamma$ 是作用在 $\Sigma$ 上的有限群，作用由全纯自同构给出，并且使得 $\Sigma$ 上 $\Gamma$-等变平坦 $\Sl(2,\C)$-联络的模空间的某个连通分支具有模空间 $\cM_{\mathfrak{g}}$ 维数的一半。则下列情形之一成立：
\begin{enumerate}[label=(\Roman*),leftmargin=2.8em]
\item $\Gamma=\Z_2$ 通过无不动点的对合作用在 $\Sigma$ 上，或
\item $\Sigma$ 是亏格 3 的超椭圆黎曼面，且 $\Gamma=\Z_2\times\Z_2$。
\end{enumerate}
在后一情形 (II) 中，其中一个 $\Z_2$-因子对应于超椭圆对合，而另一个 $\Z_2$-因子对应于一个具有 4 个不动点的对合。
\end{thm}

% [Tutor Note: p.28 定理 4.2 原文 "In the later case (II)" 应为 "In the latter case (II)"；另外定理中 "the moduli space $\cM_{\mathfrak{g}}$" 的下标是 fraktur $\mathfrak{g}$，与本讲义通篇的 $\cM_{G_\C}$ 记号不一致（疑为 $\cM_g$ 或 $\cM_{G_\C}$ 之误），照原样排。]

值得注意的是，所有带有无不动点作用的亏格 3 黎曼面必定是超椭圆的（见 [59, Proposition 21]），因此一个亏格 3 的超椭圆黎曼面连同其上一个具有 4 个不动点的额外对合，就等同于一个带有无不动点对合的亏格 3 黎曼面。

由上述可见，无不动点对合变得特别重要。特别地，对于 $\Sl(2,\C)$-Higgs 丛，[59, Proposition 23] 证明了：给定无不动点对合 $\tau$，亏格 3 黎曼面上平坦且不可约的 $\tau$-不变 $\Sl(2,\C)$-联络关于超椭圆对合是等变的。应当指出，不可约性并不是处于等变膜中的必要条件，因为黎曼面上存在平坦可约联络，它们对应于超椭圆亏格 2 商黎曼面上的不可约联络。

% [Tutor Note: p.28 原文两处笔误：(1) 段首 "furthermore" 小写（应为 "Furthermore"）；(2) "the hyperelliptic genus 2 quotient Riemann surface surface" 末尾 "surface" 重复。]

这些膜的几何可以特别地通过考察它们与 Hitchin 纤维化的相交来研究--正如上一章考虑实 Higgs 丛时所做的那样。对于 $\Sl(2,\C)$-Higgs 丛模空间中的膜，可以看到它们用 Prym 簇来表达：

\begin{thm}{4.3（[59] 定理 14）}
设 $\tau$ 是一个无不动点对合。$\tau$-等变 $(\mathrm{B},\mathrm{B},\mathrm{B})$-膜与相应 Hitchin 纤维化的一个一般纤维相交于阿贝尔簇 $\Prym(S/\tau,\Sigma/\tau)/\Z_2$，其中该纤维位于定义谱曲线 $S$ 的点之上。
\end{thm}
